% Test that explicit memoir settings and local font choices survive startup.
\PassOptionsToPackage{changecounternumbering=false}{latexbangla}
\documentclass[11pt,oneside,openany]{bookbangla}
\setbanglamainfont{Kalpurush}
\newfontfamily\quoteface{TeX Gyre Cursor}
\bookbanglasetup{quote-font={\quoteface\itshape},
  box-frame-color=black,box-background-color=black!3}
\chapterstyle{default}
\setsecnumformat{\csname the#1\endcsname\quad}
\setsecheadstyle{\raggedright\large\itshape}
\pagestyle{plain}
\pretitle{\begin{center}\Large OVERRIDE: }
\posttitle{\par\end{center}}
\setlength{\parindent}{1em}
\setlength{\parskip}{0pt}
\title{নিজের নকশা}
\author{লেখক}
\date{}
\begin{document}
\ifdim\parindent=1em\else\errmessage{bookbangla-test-failed: paragraph override}\fi
\maketitle
\mainmatter
\chapter{নিজের অধ্যায়}
\section{নিজের পরিচ্ছেদ}
\begin{chapquote}[1em]{Author}
English quote with বাংলা and \textbf{bold words}.
\end{chapquote}
\begin{exercise}[hint={A first hint with $x=2$.}]\label{ex:first}
প্রথম সমস্যা।
\end{exercise}
\begin{exercise}\label{ex:second}
দ্বিতীয় সমস্যা। আগের ইঙ্গিত এখানে থাকবে না।
\end{exercise}
\section{পরের সমস্যা}
\begin{exercise}\label{ex:third}তৃতীয় সমস্যা।\end{exercise}
\ifnum\value{exercise}=3\else\errmessage{bookbangla-test-failed: exercise continuation}\fi
\begin{example}\label{example:first}প্রথম উদাহরণ।\end{example}
\ifnum\value{example}=1\else\errmessage{bookbangla-test-failed: separate counters}\fi
\autoref{ex:first}, \autoref{ex:second}, \autoref{ex:third}, \autoref{example:first}.
\chapter{পরের অধ্যায়}
\begin{exercise}\label{ex:reset}নতুন অধ্যায়ের সমস্যা।\end{exercise}
\ifnum\value{exercise}=1\else\errmessage{bookbangla-test-failed: chapter reset}\fi
\autoref{ex:reset}.
% Authors can replace the class's numbering through memoir's counter API.
\counterwithout{example}{chapter}
\setcounter{example}{7}
\begin{example}\label{example:continuous}অষ্টম উদাহরণ।\end{example}
\counterwithin{example}{chapter}
\begin{smartlist}[label=--,leftmargin=2em,itemsep=0pt]
\item প্রথম স্তর।
\begin{smartlist}[label=+]\item দ্বিতীয় স্তর।\end{smartlist}
\end{smartlist}
\appendix
\chapter{সংযোজন}
\begin{exercise}\label{ex:appendix}সংযোজনের সমস্যা।\end{exercise}
\autoref{ex:appendix}.
\counterwithin{exercise}{section}
\section{পরিচ্ছেদভিত্তিক সমস্যা}
\begin{exercise}\label{ex:section}একটি সমস্যা।\end{exercise}
\autoref{example:continuous}, \autoref{ex:section}.
\backmatter
\chapter{শেষ কথা}
Done.
\end{document}
